\chapter{Continuous-Time Markov Chains and Queues}

In Chapter~26, our random walker hopped along an integer line at discrete clock ticks
$n = 0, 1, 2, \dots$. But real systems do not wait for a synchronized bell. A customer
arrives at a bakery at $10\!:\!04\!:\!17$, not at ``step $7$''; a packet hits a router or a
machine breaks down at an unpredictable continuous instant.

This chapter rebuilds the Markov idea in \emph{continuous time}. Instead of step-by-step jump
probabilities, a single table of instantaneous flow rates called the \textbf{generator matrix}
captures all the dynamics. From it, we derive the steady-state behavior of \textbf{queues} ---
waiting lines that govern computer networks, hospitals, and server pools. We will meet the
fundamental single-server M/M/1 queue, watch waiting customers settle into a geometric equilibrium,
and prove \textbf{Little's law} ($L = \lambda W$).

The recipe is simple: take an ordinary Markov chain and glue exponential, memoryless waiting
times onto each state. This chapter builds that gluing from scratch.

\chapterline

\section{From Discrete Ticks to Continuous Flow}
\label{ch27:intro}

\begin{intuitionbox}
Picture a tiny machine that is either \textbf{working} or \textbf{broken}. In discrete time we would
say: each minute, if it is working it breaks with probability $0.01$; if it is broken it gets fixed
with probability $0.05$. That is a fine model, but it has an arbitrary unit baked in --- the minute.
Why not the second? The half-minute? If we shrink the time step, we must shrink the probabilities to
match, and in the limit the natural object is not a \emph{probability per step} but a \emph{rate per
unit time}. ``This machine breaks at rate $0.01$ per minute'' makes sense at every time scale at
once. Rates, not per-step probabilities, are the native language of continuous time.
\end{intuitionbox}

Making the limit explicit shows exactly where the exponential distribution comes
from. Suppose a working machine breaks, in each tiny interval of length $h$, with probability
approximately $\lambda h$ (proportional to the length of the interval --- twice the time, twice the
chance, for small $h$). Cut the first $t$ units of time into $n = t/h$ slices. The probability of
\emph{surviving} all of them, if breakdowns in disjoint slices are independent, is
\[
  \P(\text{still working at time } t)
  \;\approx\; (1 - \lambda h)^{n}
  \;=\; \left(1 - \frac{\lambda t}{n}\right)^{n}
  \;\xrightarrow[n \to \infty]{}\; e^{-\lambda t}.
\]
That is the survival function of an $\Expo(\lambda)$ random variable. The lifetime of the machine
before its first breakdown is exponential, and the parameter $\lambda$ is the exact \emph{rate} we
started with --- the engine of the entire chapter.

\begin{whybox}
Why must the holding time be exponential, and nothing else? Because we insisted on the
\textbf{Markov property in continuous time}: the future depends on the present state alone, not on
how long we have already been sitting in it. If the time spent in a state had any other distribution
--- say, ``a machine that has run a long time is more likely to break soon'' --- then knowing the
\emph{age} of the current sojourn would tell you something about the future, violating
memorylessness. The exponential distribution is the \emph{unique} continuous distribution with no
memory:
\[
  \P(T > s + t \given T > s) = \P(T > t)
  \quad\text{for all } s, t \ge 0.
\]
So ``continuous-time Markov'' forces ``exponential holding times.'' The two are the same demand
spoken in two languages.
\end{whybox}

\begin{recoverybox}
\textbf{Pause and think:} Suppose a lightbulb has an exponential lifetime with failure rate
$\lambda = 0.001$ per hour (mean lifetime $1000$ hours). You enter a room and find the bulb already
on. How long has it been glowing? You have no way of knowing. But how much longer do you expect it
to last from this very second? By memorylessness, still exactly $1/\lambda = 1000$ hours. Its elapsed
history gives you zero information about its remaining life.
\end{recoverybox}

\subsection{The two-part recipe}

Here is the cleanest way to picture a continuous-time Markov chain (CTMC), and we will return to it
constantly. The process sits in some state $i$. Two things happen:

\begin{enumerate}
  \item \textbf{It waits.} The time it spends in state $i$ (the \emph{holding time} or
  \emph{sojourn time}) is exponential with some rate $q_i$ that depends only on $i$. On average it
  waits $1/q_i$.
  \item \textbf{It jumps.} When the holding time expires, it jumps to a new state $j \ne i$ chosen
  by a probability $P_{ij}$ (an ordinary discrete transition, with $\sum_{j \ne i} P_{ij} = 1$).
\end{enumerate}

Then it does the same thing from $j$, and so on. The embedded sequence of states it visits ---
ignoring how long it stayed --- is just an ordinary discrete-time Markov chain, called the
\emph{jump chain} or \emph{embedded chain}. The continuous-time part lives entirely in the
exponential waiting times glued onto each visit.

\begin{happeningbox}
A CTMC combines a discrete jump chain with exponential holding times. Without the random sojourn times,
you have a familiar jump chain on the same states. Adding them back yields a process that
lives in real time. The generator, forward equations, and balance conditions keep these two
layers (where it jumps and how long it stays) consistent with one another.
\end{happeningbox}

\section{Rates and the Generator Matrix}
\label{ch27:generator}

The two-part recipe (a rate $q_i$ to leave each state, plus jump probabilities $P_{ij}$) is intuitive
but clumsy to compute with. A single object packages both and makes the algebra manageable:
the \emph{generator matrix} $Q$, also called the \emph{rate matrix} or \emph{infinitesimal
generator}.

\begin{intuitionbox}
The transition matrix of a discrete chain answers ``where will I be \emph{next step}?'' But in
continuous time there is no ``next step'' --- time flows. So instead the generator answers a
derivative question: ``\emph{how fast} is probability flowing out of where I am, and into each other
state, \emph{right now}?'' It is the time-derivative of transition probability at $t = 0$. Rates,
not probabilities --- and that shift of viewpoint is the whole definition.
\end{intuitionbox}

\begin{definition}[Transition rates and generator]
\label{ch27:def-generator}
For a CTMC, define the \emph{off-diagonal rate} from state $i$ to state $j \ne i$ as
\[
  q_{ij} \;=\; \lim_{h \to 0^+} \frac{\P(X_{t+h} = j \given X_t = i)}{h},
\]
the instantaneous rate of jumping $i \to j$. Define the \emph{total exit rate} from $i$ as
$q_i = \sum_{j \ne i} q_{ij}$, and set the diagonal entry $q_{ii} = -q_i$. The matrix
$Q = (q_{ij})$ is the \emph{generator}. Its defining properties are
\[
  q_{ij} \ge 0 \ (i \ne j), \qquad q_{ii} = -\sum_{j \ne i} q_{ij},
  \qquad\text{so every row of } Q \text{ sums to } 0.
\]
\end{definition}

The two viewpoints connect cleanly. The holding-time rate in state $i$ is the total exit rate
$q_i = \sum_{j \ne i} q_{ij}$, so the sojourn time is $\Expo(q_i)$. And when the chain does leave,
it goes to $j$ with probability
\[
  P_{ij} \;=\; \frac{q_{ij}}{q_i},
\]
the fraction of the total outflow rate that heads toward $j$. So $Q$ contains \emph{both} parts of
the recipe: the row sums of off-diagonals give the holding rates, and their proportions give the
jump probabilities.

\begin{walkbox}
Suppose state $1$ has $q_{12} = 3$ and $q_{13} = 1$ (rates per hour). How long, on average, does the
chain stay in state $1$, and where does it go when it leaves?

The total exit rate is $q_1 = 3 + 1 = 4$ per hour, so the mean holding time is $1/4$ hour $= 15$
minutes. When it leaves, it goes to state $2$ with probability $3/4$ and to state $3$ with
probability $1/4$. The faster a particular exit ``ticks,'' the more likely it is to be the one that
fires --- a race between exponential alarm clocks, won most often by the fastest.
\end{walkbox}

\begin{connectionbox}
The race-of-exponentials picture is worth keeping. Imagine that from state $i$, every possible
target $j$ has its own independent exponential alarm clock with rate $q_{ij}$. The chain jumps when
the \emph{first} clock rings. Two facts from the Poisson/exponential chapter make this work: (i) the
minimum of independent exponentials with rates $q_{ij}$ is itself exponential with rate
$\sum_j q_{ij} = q_i$; and (ii) the probability that clock $j$ is the winner is
$q_{ij}/q_i$, independent of \emph{when} the winner rings. So ``waiting $\Expo(q_i)$ then jumping by
$P_{ij}$'' and ``racing independent clocks'' are the same process. The generator is just the table of
all the clock rates.
\end{connectionbox}

\subsection{From the generator to transition probabilities}

Let $P_{ij}(t) = \P(X_{t+s} = j \given X_s = i)$ be the probability that the chain is in state $j$
after an elapsed time $t$, given it began in $i$. Gather these into a transition matrix $P(t)$.
At time zero the chain has had no time to move, so $P(0) = I$ (the identity matrix). For a tiny time
increment $h > 0$, our definition of the generator $Q$ in Definition~\ref{ch27:def-generator} says
\[
  P(h) \;\approx\; I + Q h,
\]
since each off-diagonal transition probability is $\approx q_{ij}h$ and the probability of staying in
state $i$ is $\approx 1 - q_i h = 1 + q_{ii}h$.

Now compose an elapsed time $t$ with an extra tiny interval $h$ in two different orders:
\begin{itemize}
  \item \textbf{Looking at the end (Forward):} Run the chain for time $t$, then watch the final
  sliver $h$. By the Markov property, $P(t+h) = P(t)P(h) \approx P(t)(I + Qh) = P(t) + P(t)Qh$.
  Subtracting $P(t)$ and dividing by $h$ yields the \textbf{Kolmogorov forward equation}:
  \[
    P'(t) \;=\; P(t)\,Q.
  \]
  \item \textbf{Looking at the start (Backward):} Watch the first sliver $h$, then let the chain
  evolve for the remaining time $t$. Here $P(h+t) = P(h)P(t) \approx (I + Qh)P(t) = P(t) + QP(t)h$.
  This yields the \textbf{Kolmogorov backward equation}:
  \[
    P'(t) \;=\; Q\,P(t).
  \]
\end{itemize}
Both equations describe the exact same rate of change. Forward watches the finish line; backward
watches the starting gate. Because $Q$ and $P(t)$ commute ($Qe^{Qt} = e^{Qt}Q$), both yield the
same matrix-exponential solution:
\[
  \boxed{\,P(t) = e^{Qt} = \sum_{k=0}^{\infty} \frac{(Qt)^k}{k!}\,.}
\]

\begin{connectionbox}
Compare with the scalar story. The ordinary differential equation $p'(t) = -\lambda p(t)$ has
solution $p(t) = p(0)e^{-\lambda t}$ --- exponential decay, the survival curve from
Section~\ref{ch27:intro}. The CTMC obeys exactly the same law, only now $\lambda$ is replaced by the
matrix $-Q$ and the scalar exponential becomes the \emph{matrix exponential}. If you have met
matrix exponentials in a linear algebra course (eigenvalues, diagonalization, $e^{Qt}$ acting
through eigenvectors), this is their direct role in probability: the generator's
eigenvalues control how fast the chain forgets its start and relaxes to equilibrium.
\end{connectionbox}

\begin{example}[A two-state machine]
\label{ch27:ex-twostate}
Let state $0$ = ``broken,'' state $1$ = ``working.'' A broken machine is repaired at rate $\mu$, a
working machine breaks at rate $\lambda$. The generator is
\[
  Q = \begin{pmatrix} -\mu & \mu \\ \lambda & -\lambda \end{pmatrix}.
\]
Each row sums to zero, as it must. The holding time when broken is $\Expo(\mu)$ (mean repair time
$1/\mu$); when working it is $\Expo(\lambda)$ (mean uptime $1/\lambda$). Solving $\pi Q = 0$ with
$\pi_0 + \pi_1 = 1$ (we justify this stationarity equation in Section~\ref{ch27:steady}) gives the
long-run fraction of time working,
\[
  \pi_1 = \frac{\mu}{\lambda + \mu},
  \qquad \pi_0 = \frac{\lambda}{\lambda + \mu}.
\]
If repairs are fast ($\mu$ large) relative to failures ($\lambda$ small), the machine is up almost
all the time --- the ratio $\mu/(\lambda+\mu)$ is exactly the \emph{availability} an engineer would
quote. Notice it depends only on the rates, not on the precise distribution shapes beyond their
means; that robustness is a recurring theme.
\end{example}

\begin{figure}[H]
\centering
\begin{tikzpicture}[->,>=Stealth,node distance=3.6cm,on grid,auto]
  \node[draw,circle,minimum size=1.1cm,fill=densityred!12] (b) {broken};
  \node[draw,circle,minimum size=1.1cm,fill=targetgreen!15,right=of b] (w) {working};
  \draw[sourceblue,bend left=28] (b) to node[above] {repair rate $\mu$} (w);
  \draw[densityred,bend left=28] (w) to node[below] {failure rate $\lambda$} (b);
\end{tikzpicture}
\caption{The two-state machine of Example~\ref{ch27:ex-twostate}. The chain hops between broken and
working; the long-run fraction of time working is $\mu/(\lambda+\mu)$, larger when repairs (rate
$\mu$) outpace failures (rate $\lambda$). This is the smallest interesting CTMC and the building block
of reliability models.}
\label{ch27:fig-twostate}
\end{figure}

\section{Birth-Death Processes: Counting Up and Down}
\label{ch27:birthdeath}

A vast and useful family of CTMCs lives on the states $\{0, 1, 2, \dots\}$ --- think of them as a
\emph{population size}, a \emph{number of customers}, or a \emph{queue length} --- and can only move
to a neighbor: up by one (a \emph{birth}, or an arrival) or down by one (a \emph{death}, or a
departure). These are \emph{birth-death processes}, and they are foundational models in queueing theory.

\begin{definition}[Birth-death process]
\label{ch27:def-bd}
A birth-death process on $\{0,1,2,\dots\}$ has \emph{birth rates} $\lambda_n \ge 0$ (rate of going
$n \to n+1$) and \emph{death rates} $\mu_n \ge 0$ (rate of going $n \to n-1$, with $\mu_0 = 0$ since
the population cannot go negative). All other rates are zero. The generator is tridiagonal:
\[
  q_{n,n+1} = \lambda_n, \qquad q_{n,n-1} = \mu_n,
  \qquad q_{n,n} = -(\lambda_n + \mu_n).
\]
\end{definition}

\begin{figure}[H]
\centering
\begin{tikzpicture}[->,>=Stealth,node distance=2.3cm,on grid,auto,
  mystate/.style={circle,draw=gray!80,thick,fill=gray!8,minimum size=0.95cm,font=\small\bfseries}]
  \node[mystate] (0) {$0$};
  \node[mystate,right=of 0] (1) {$1$};
  \node[mystate,right=of 1] (2) {$2$};
  \node[mystate,right=of 2] (3) {$3$};
  \node[right=1.8cm of 3, font=\small] (4) {$\cdots$};

  \draw[sourceblue,thick,bend left=32] (0) to node[above, font=\footnotesize] {$\lambda_0$} (1);
  \draw[sourceblue,thick,bend left=32] (1) to node[above, font=\footnotesize, pos=0.25] {$\lambda_1$} (2);
  \draw[sourceblue,thick,bend left=32] (2) to node[above, font=\footnotesize] {$\lambda_2$} (3);
  \draw[sourceblue,thick,bend left=32] (3) to node[above, font=\footnotesize] {$\lambda_3$} (4);

  \draw[densityred,thick,bend left=32] (1) to node[below, font=\footnotesize] {$\mu_1$} (0);
  \draw[densityred,thick,bend left=32] (2) to node[below, font=\footnotesize, pos=0.25] {$\mu_2$} (1);
  \draw[densityred,thick,bend left=32] (3) to node[below, font=\footnotesize] {$\mu_3$} (2);
  \draw[densityred,thick,bend left=32] (4) to node[below, font=\footnotesize] {$\mu_4$} (3);

  % Cut boundary between states 1 and 2
  \draw[dashed, structurepurple, line width=1.2pt] (3.45, -1.3) -- (3.45, 1.3)
    node[above, font=\scriptsize\bfseries, structurepurple, align=center] {Cut boundary\\ $\pi_1\lambda_1 = \pi_2\mu_2$};
\end{tikzpicture}
\caption{A birth-death process with detailed balance cut. Blue arrows are births (arrivals) at rates $\lambda_n$; red arrows are deaths (departures) at rates $\mu_n$. In steady state, probability mass cannot pool; the flow across any vertical cut (purple dashed line) must balance: $\pi_n \lambda_n = \pi_{n+1} \mu_{n+1}$.}
\label{ch27:fig-bd}
\end{figure}

\subsection{Detailed balance and the product formula}

What is the long-run distribution $\pi_n$ over states? Here the neighbor-only structure pays off
spectacularly. In equilibrium, consider the ``cut'' between states $\{0, \dots, n\}$ and
$\{n+1, n+2, \dots\}$. The only way to cross it upward is the birth $n \to n+1$; the only way to
cross it downward is the death $n+1 \to n$. In steady state the flow across any such cut must
balance --- whatever goes up must come back down, or probability would pile up on one side forever.
This gives the \emph{detailed balance} equations:
\[
  \pi_n \lambda_n \;=\; \pi_{n+1}\,\mu_{n+1}
  \qquad\text{for all } n \ge 0.
\]

\begin{whybox}
Why does flow have to balance across every cut? Imagine a wall between ``$\le n$'' and ``$\ge n+1$.''
Probability mass leaks across it rightward at rate $\pi_n \lambda_n$ (you must be in state $n$, and
births fire at rate $\lambda_n$) and leftward at rate $\pi_{n+1}\mu_{n+1}$. If these were unequal,
the total probability on one side would be steadily growing --- but in a \emph{stationary}
distribution nothing changes with time, by definition. So the net flow across the wall is zero. For
birth-death chains this cut argument gives the equilibrium directly, no matrix algebra required. (The
same trick works for any chain whose graph is a tree, because every edge is a cut.)
\end{whybox}

\begin{recoverybox}
\textbf{Pause and think:} What would happen physically if cut flow were unbalanced, say
$\pi_0 \lambda_0 > \pi_1 \mu_1$? In every second, more probability mass would cross into states
$\{1, 2, \dots\}$ than returned down to state $0$. Probability would steadily drain out of state $0$
and pool in higher states. But in a \emph{stationary} distribution, by definition, the probability
$\pi_0$ cannot change over time! The puddle can neither grow nor drain. So upward flow must exactly
match downward flow across every cut boundary.
\end{recoverybox}

Solving the recursion $\pi_{n+1} = \pi_n \cdot \frac{\lambda_n}{\mu_{n+1}}$ step by step reveals the pattern:
\[
  \pi_1 = \pi_0 \frac{\lambda_0}{\mu_1}, \qquad
  \pi_2 = \pi_1 \frac{\lambda_1}{\mu_2} = \pi_0 \frac{\lambda_0 \lambda_1}{\mu_1 \mu_2}, \qquad
  \pi_3 = \pi_2 \frac{\lambda_2}{\mu_3} = \pi_0 \frac{\lambda_0 \lambda_1 \lambda_2}{\mu_1 \mu_2 \mu_3},
\]
and by induction, this telescopes into the general \textbf{product formula}:
\[
  \boxed{\,\pi_n = \pi_0 \prod_{k=0}^{n-1} \frac{\lambda_k}{\mu_{k+1}}\,,}
  \qquad
  \pi_0 = \left(1 + \sum_{n=1}^{\infty} \prod_{k=0}^{n-1}\frac{\lambda_k}{\mu_{k+1}}\right)^{-1}.
\]
The normalizing constant $\pi_0$ just makes the $\pi_n$ sum to $1$; it exists (the chain has a steady
state) precisely when that infinite sum converges. Hold onto this product formula --- the M/M/1
queue is the simplest case, where all the birth rates are equal and all the death rates are equal,
and the product collapses into a geometric distribution.

\section{The M/M/1 Queue}
\label{ch27:mm1}

Now the centerpiece. A single server --- one barista, one CPU, one checkout lane --- handles
customers one at a time. We make the cleanest possible assumptions, encoded in the cryptic name
\textbf{M/M/1}:

\begin{itemize}
  \item \textbf{First M} (\emph{Markovian arrivals}): customers arrive according to a Poisson process
  with rate $\lambda$. Equivalently, the gaps between arrivals are i.i.d.\ $\Expo(\lambda)$.
  \item \textbf{Second M} (\emph{Markovian service}): each customer's service time is $\Expo(\mu)$,
  independent of everything else.
  \item \textbf{1}: a single server.
\end{itemize}

The ``M'' stands for \emph{memoryless} (sometimes said to honor Markov), and that memorylessness is
exactly what makes the queue length a Markov chain. Let $X_t$ be the number of customers in the
system at time $t$ (those waiting plus the one in service). Because arrivals come at constant rate
$\lambda$ and service completes at constant rate $\mu$ whenever someone is being served, $X_t$ is a
birth-death process with
\[
  \lambda_n = \lambda \ (\text{for all } n \ge 0),
  \qquad \mu_n = \mu \ (\text{for all } n \ge 1).
\]
Figure~\ref{ch27:fig-samplepath} shows a typical realization of $X_t$: a staircase that jumps up at
arrivals and down at service completions.

\begin{historybox}
The mathematics of queues was born of a concrete engineering problem: how many telephone lines
does a city need? In 1909, the Danish mathematician and engineer \textbf{Agner Krarup Erlang}
(1878--1929), working for the Copenhagen Telephone Company, set out to size switchboards and
underground copper cables connecting Copenhagen's exchanges. If the company installed too few lines,
callers received busy signals; if it laid too many, copper sat idle. Erlang sought an exact formula
for the probability that a caller finds every line occupied.

Decades before Andrei Kolmogorov formalized probability (1931), Erlang balanced the rate of probability
flowing into each state against the rate flowing out. His insight yielded two foundational formulas:
\emph{Erlang B} (for loss systems, where a caller finding all lines busy is blocked) and \emph{Erlang C}
(for delay systems, where callers wait in line for the next available operator). In 1946, the
dimensionless unit of telecommunications traffic, the \emph{erlang}, was named in his honor.
Queueing theory was soon extended to non-exponential service times by Felix Pollaczek (1930) and
Aleksandr Khinchin (1932) in the celebrated M/G/1 queue.

Yet major questions remain open. In the 1990s, Walter Leland and colleagues showed that real internet
packet arrivals are \emph{self-similar} --- clustering into heavy-tailed bursts across time scales,
violating Poisson assumptions. Today, predicting delay in multi-server queues with general service
distributions ($G/G/c$) and designing optimal load-balancing policies for large-scale cloud clusters
remain active areas of research.
\end{historybox}

\subsection{The traffic intensity and steady state}

Define the \emph{traffic intensity} (or utilization)
\[
  \rho \;=\; \frac{\lambda}{\mu},
\]
the ratio of how fast work arrives to how fast it can be done. This single dimensionless number
controls everything. Plug the constant rates into the birth-death product formula:
\[
  \pi_n = \pi_0 \prod_{k=0}^{n-1}\frac{\lambda}{\mu} = \pi_0\, \rho^{\,n}.
\]
For the $\pi_n$ to sum to one we need $\sum_n \rho^n < \infty$, i.e.\ $\rho < 1$: \emph{the server
must be faster than the arrivals.} When $\rho < 1$ the geometric series sums to $1/(1-\rho)$, so
$\pi_0 = 1 - \rho$ and
\[
  \boxed{\,\pi_n = (1-\rho)\,\rho^{\,n}, \quad n = 0, 1, 2, \dots\,}
\]
The number in the system has a \emph{geometric distribution}. The probability the server is idle is
$\pi_0 = 1-\rho$; the long-run fraction of time it is busy is $\rho$ --- which is why $\rho$ is also
called the utilization.

\begin{figure}[H]
\centering
\begin{tikzpicture}
\begin{axis}[
  width=12cm, height=6.2cm,
  ybar, bar width=7pt,
  xlabel={number in system $n$}, ylabel={steady-state probability $\pi_n$},
  xtick={0,1,2,...,12}, ymin=0, ymax=0.62,
  title style={font=\small},
  title={Steady-state M/M/1 queue length distribution $\pi_n = (1-\rho)\rho^n$},
  legend style={at={(0.97,0.97)},anchor=north east,font=\small},
  enlarge x limits=0.04, tick label style={font=\small},
]
\addplot[draw=sourceblue,fill=sourceblue!55] coordinates {
  (0,0.5)(1,0.25)(2,0.125)(3,0.0625)(4,0.03125)(5,0.015625)(6,0.0078125)
  (7,0.00390625)(8,0.001953)(9,0.000977)(10,0.000488)(11,0.000244)(12,0.000122)};
\addlegendentry{$\rho = 0.5$}
\addplot[draw=densityred,fill=densityred!45] coordinates {
  (0,0.2)(1,0.16)(2,0.128)(3,0.1024)(4,0.08192)(5,0.065536)(6,0.052429)
  (7,0.041943)(8,0.033554)(9,0.026844)(10,0.021475)(11,0.01718)(12,0.013744)};
\addlegendentry{$\rho = 0.8$}
\end{axis}
\end{tikzpicture}
\caption{The steady-state queue length of an M/M/1 system is geometric, $\pi_n = (1-\rho)\rho^n$. At
low load ($\rho = 0.5$, blue) the system is usually nearly empty. At high load ($\rho = 0.8$, red)
the distribution has a long tail --- long lines become common. As $\rho \to 1$ the tail stretches
without bound.}
\label{ch27:fig-mm1}
\end{figure}

\begin{figure}[H]
\centering
\begin{tikzpicture}
\begin{axis}[
  width=12.5cm, height=5.4cm,
  xlabel={time $t$}, ylabel={number in system $X_t$},
  ymin=0, ymax=5, xmin=0, xmax=20,
  ytick={0,1,2,3,4}, xtick={0,4,8,12,16,20},
  title style={font=\small},
  title={Sample path of the M/M/1 queue: continuous-time jump process},
  tick label style={font=\small}, const plot,
]
\addplot[sourceblue,very thick,const plot] coordinates {
  (0,0)(1.2,1)(2.0,2)(3.1,1)(4.4,2)(5.0,3)(6.3,2)(7.1,1)(8.0,0)
  (9.5,1)(10.2,2)(11.0,1)(12.4,2)(13.1,3)(13.9,4)(15.0,3)(16.2,2)
  (17.0,1)(18.3,2)(19.1,1)(20,1)};
\end{axis}
\end{tikzpicture}
\caption{A sample path of the M/M/1 queue length $X_t$: a right-continuous staircase that steps up by
one at each arrival and down by one at each service completion. Holding times in each level are
$\Expo$ (rate $\lambda$ up, rate $\mu$ down once busy); the steady-state histogram of the heights is
the geometric law of Figure~\ref{ch27:fig-mm1}.}
\label{ch27:fig-samplepath}
\end{figure}

\subsection{Performance measures}

From the geometric distribution we read off how the system performs. The mean number in the system
is the mean of a geometric:
\[
  L \;=\; \E[X] \;=\; \sum_{n=0}^{\infty} n(1-\rho)\rho^n \;=\; \frac{\rho}{1-\rho}.
\]
This is the formula that should keep every system designer awake at night. Watch what happens as the
server gets busy:
\[
  \rho = 0.5 \Rightarrow L = 1, \qquad
  \rho = 0.9 \Rightarrow L = 9, \qquad
  \rho = 0.99 \Rightarrow L = 99.
\]

\begin{figure}[H]
\centering
\begin{tikzpicture}
\begin{axis}[
  width=12cm, height=6cm,
  xlabel={utilization $\rho = \lambda/\mu$}, ylabel={mean number in system $L$},
  domain=0:0.97, samples=120, xmin=0, xmax=1, ymin=0, ymax=35,
  title style={font=\small},
  title={Hyperbolic backlog explosion: mean queue length $L = \rho/(1-\rho)$},
  thick, tick label style={font=\small},
]
\addplot[densityred,very thick] {x/(1-x)};
\draw[gray,dashed] (axis cs:0.9,0) -- (axis cs:0.9,9) -- (axis cs:0,9);
\node[above left,font=\scriptsize,gray!80!black] at (axis cs:0.9,9) {$(0.9,\,9)$};
\node[anchor=west,font=\small] at (axis cs:0.55,28) {$L = \dfrac{\rho}{1-\rho}$};
\end{axis}
\end{tikzpicture}
\caption{The mean number in an M/M/1 queue explodes as utilization approaches $1$. The dashed lines
mark $\rho = 0.9$, where $L = 9$. The blow-up is hyperbolic: pushing utilization from $90\%$ to
$99\%$ multiplies the average backlog tenfold. This nonlinearity is the central practical lesson of
queueing theory.}
\label{ch27:fig-L}
\end{figure}

\begin{trapbox}
The deadly intuition is that a server running at $99\%$ utilization is ``almost the same'' as one at
$90\%$ --- after all, it is only $9$ percentage points busier, and you are squeezing nearly all the
work out of your expensive hardware. The mathematics says otherwise: the queue is ten times longer.
The reason is \emph{variability}. Even though work arrives at rate $\lambda$ and is served at rate
$\mu > \lambda$ \emph{on average}, the randomness in arrivals and service means temporary pile-ups
happen, and the fuller the system, the longer it takes to drain them. Near $\rho = 1$ there is almost
no idle capacity to absorb a burst, so backlogs linger. \textbf{Provision for the variance, not just
the mean.} A system sized so that average load equals average capacity ($\rho = 1$) is not balanced
--- it is unstable, with an infinitely long expected queue.
\end{trapbox}

\begin{paradoxbox}
Two engineers report on a server farm. One says ``utilization is $50\%$ --- we are wasting half our
capacity, buy fewer machines.'' The other says ``utilization is $99\%$ --- we are perfectly
efficient, our machines are never idle.'' Which farm would you rather be a customer of? The
high-utilization farm has, on average, $99$ jobs backed up versus the relaxed farm's $1$. The
``efficient'' system is the miserable one to wait in. The paradox dissolves once you separate two
goals that pull in opposite directions: \emph{keeping the server busy} (good for the owner's
wallet) and \emph{keeping the wait short} (good for the customer). High utilization optimizes the
first and destroys the second. There is no setting that does both, because $L = \rho/(1-\rho)$ ties
them together. Real capacity planning is choosing where on this curve to sit --- and the right answer
is almost never near the top.
\end{paradoxbox}

From the backlog $L$, we can immediately deduce the customer's average waiting experience.
Let $W$ be the mean total time a customer spends in the system (waiting in line plus being
served). How long does an arriving customer expect to stay?
Imagine stepping into the shop: by the memoryless property of Poisson arrivals (PASTA, discussed
below), an arrival catches the queue in its typical equilibrium, seeing on average $L$ customers
ahead.

Because service times are memoryless $\Expo(\mu)$, the customer currently at the counter requires an
average of $1/\mu$ time units to finish, regardless of how long they have already been there. Each of
the remaining $L-1$ customers ahead takes an independent mean time $1/\mu$, and the new arrival's
own service takes $1/\mu$. That is a total of $L + 1$ service completions, each of average duration
$1/\mu$. Therefore:
\[
  W \;=\; \frac{L + 1}{\mu}
     \;=\; \frac{\frac{\rho}{1-\rho} + 1}{\mu}
     \;=\; \frac{\frac{1}{1-\rho}}{\mu}
     \;=\; \frac{1}{\mu(1 - \lambda/\mu)}
     \;=\; \frac{1}{\mu - \lambda}.
\]
Subtracting the customer's own mean service time $1/\mu$ yields the average time spent purely waiting
in line,
\[
  W_q \;=\; W - \frac{1}{\mu}
      \;=\; \frac{1}{\mu - \lambda} - \frac{1}{\mu}
      \;=\; \frac{\mu - (\mu - \lambda)}{\mu(\mu - \lambda)}
      \;=\; \frac{\lambda}{\mu(\mu - \lambda)}
      \;=\; \frac{\rho}{\mu - \lambda}.
\]
And the average number of customers waiting in the line itself (excluding the one being served) is
the total backlog minus the fraction of time the server is occupied:
\[
  L_q \;=\; L - \rho \;=\; \frac{\rho}{1-\rho} - \rho \;=\; \frac{\rho^2}{1-\rho}.
\]

\begin{example}[Sizing a help desk]
\label{ch27:ex-helpdesk}
A help desk receives $\lambda = 18$ tickets per hour. A single agent resolves tickets at
$\mu = 20$ per hour. Then $\rho = 18/20 = 0.9$. The mean number of tickets in the system is
$L = 0.9/0.1 = 9$, and the mean time from submission to resolution is
$W = 1/(20 - 18) = 0.5$ hours $= 30$ minutes. Management wants to halve the wait. Doubling staff is
the obvious move, but notice the leverage in $\mu$: if a better tool lets the agent resolve
$\mu = 24$ per hour, then $\rho = 0.75$, $L = 3$, and $W = 1/6$ hour $= 10$ minutes --- a $20\%$
speedup in service time cut the wait by two-thirds, because it pushed $\rho$ away from the cliff edge
near $1$. Small improvements in service rate buy outsized improvements in delay when you are highly
loaded.
\end{example}

\begin{labbox}
Before reading on, predict the shape of the answer to this: in the M/M/1 system, what fraction of
\emph{arriving customers} find the server idle (and so are served immediately)?

\textbf{Answer:} By the celebrated PASTA principle (``Poisson Arrivals See Time Averages''), an
arrival arriving purely at random does not anticipate busy or idle periods --- it takes an unbiased
statistical snapshot of the system in steady state. So the fraction of arrivals finding an idle
server is simply the stationary probability of state zero: $\pi_0 = 1 - \rho$.

With $\rho = 0.9$, only $10\%$ of customers walk straight up to an idle agent; the other $90\%$ must
wait. This is why a service that ``feels horribly congested'' can still have a modest average
utilization --- customers experience the system when it is full, not when it is empty.
\end{labbox}

\section{Little's Law}
\label{ch27:little}

We keep relating $L$ (how many are in the system) to $W$ (how long each one stays). The bridge
between them is one of the most general results in applied probability, and it requires almost no
assumptions.

\begin{theorem}[Little's law]
\label{ch27:thm-little}
For any stable queueing system in steady state, with $\lambda$ the long-run average arrival rate,
$L$ the long-run average number of customers in the system, and $W$ the long-run average time a
customer spends in the system,
\[
  \boxed{\,L = \lambda\, W.\,}
\]
\end{theorem}

\begin{intuitionbox}
Little's law is just a statement about \emph{accounting}, and the cleanest way to see it is to
imagine each customer pays \$1 per minute for every minute they sit in the system. There are two ways
to add up the total money collected per unit time. \textbf{Way one (look at the room):} at any
instant there are on average $L$ customers present, each paying \$1/minute, so money pours in at rate
$L$ dollars per minute. \textbf{Way two (look at the door):} customers arrive at rate $\lambda$ per
minute, and each one eventually pays $W$ dollars total (\$1/minute for their $W$ minutes inside), so
money is committed at rate $\lambda W$ dollars per minute. Both count the same dollars. Therefore
$L = \lambda W$. No distributions, no independence, no Markov property --- just two ways of
tallying the same revenue.
\end{intuitionbox}

\begin{whybox}
Now to make the accounting argument precise. Over a long horizon $[0, T]$, let $A(T)$ be the number
of arrivals and let customer $i$ spend time $W_i$ in the system. The \emph{area under the
occupancy curve} --- the total customer-minutes accumulated --- can be computed two ways. Summing
over time, it is $\int_0^T X_t\,\dd t$, which over a long horizon averages to $L \cdot T$. Summing
over customers, it is $\sum_{i=1}^{A(T)} W_i$ (each customer contributes a horizontal sliver of
width $W_i$), which equals $A(T) \cdot \overline{W}$ where $\overline{W}$ is the average sojourn.
Equate and divide by $T$:
\[
  L \;=\; \frac{1}{T}\int_0^T X_t\,\dd t
       \;=\; \frac{A(T)}{T}\cdot \overline{W}
       \;\xrightarrow[T \to \infty]{}\; \lambda \cdot W,
\]
since $A(T)/T \to \lambda$. The only thing we needed was that the system is stable (everybody
eventually leaves, so the area is finite per unit time) and that the long-run averages exist. That is
why Little's law holds for the M/M/1 queue, the M/M/$c$ queue, a grocery store, a hospital ward, or
the number of half-finished features in a software team's pipeline.
\end{whybox}

\begin{happeningbox}
Little's law is really a theorem about \emph{areas}. Draw each customer's visit as a horizontal bar
of length $W_i$. The total area accumulated over a long window $[0, T]$ represents the total
``customer-minutes'' experienced.

That identical area can be sliced two ways (Figure~\ref{ch27:fig-little-area}):
\begin{itemize}
  \item \textbf{Vertically (by time):} At each instant $t$, the vertical height is the number of
  customers present, $X_t$. Integrating over time yields $\int_0^T X_t\,\dd t \approx L \cdot T$.
  \item \textbf{Horizontally (by customer):} Each customer contributes a bar of length $W_i$.
  Summing across all $A(T) \approx \lambda T$ arrivals gives $\sum W_i \approx \lambda T \cdot W$.
\end{itemize}
Two ways to compute one area --- that is the entire content. Equating $L \cdot T = \lambda T \cdot W$
and dividing by $T$ leaves $L = \lambda W$. This is the continuous analogue of the classic Fubini
double-counting trick from discrete combinatorics. Once you see it as an area, the law's universality
stops being mysterious.
\end{happeningbox}

\begin{figure}[H]
\centering
\begin{tikzpicture}[>=Stealth,scale=0.92]
  \draw[->, thick] (0,0) -- (11.5,0) node[right, font=\small] {time $t$};
  \draw[->, thick] (0,0) -- (0,4.2) node[above, font=\small] {occupancy / customers};

  \draw[dashed, gray!70] (10,0) -- (10,3.8);
  \node[below, font=\small] at (10,0) {$T$};

  \draw[sourceblue, very thick] (0.6, 0.7) -- (3.6, 0.7) node[midway, above, font=\scriptsize, sourceblue] {Customer 1 ($W_1$)};
  \fill[sourceblue] (0.6,0.7) circle (2pt);
  \fill[sourceblue] (3.6,0.7) circle (2pt);

  \draw[targetgreen, very thick] (1.8, 1.5) -- (6.2, 1.5) node[pos=0.25, above, font=\scriptsize, targetgreen] {Customer 2 ($W_2$)};
  \fill[targetgreen] (1.8,1.5) circle (2pt);
  \fill[targetgreen] (6.2,1.5) circle (2pt);

  \draw[quotientorange, very thick] (2.7, 2.3) -- (5.1, 2.3) node[pos=0.22, above, font=\scriptsize, quotientorange] {Customer 3 ($W_3$)};
  \fill[quotientorange] (2.7,2.3) circle (2pt);
  \fill[quotientorange] (5.1,2.3) circle (2pt);

  \draw[structurepurple, very thick] (4.3, 3.1) -- (8.3, 3.1) node[pos=0.72, above, font=\scriptsize, structurepurple] {Customer 4 ($W_4$)};
  \fill[structurepurple] (4.3,3.1) circle (2pt);
  \fill[structurepurple] (8.3,3.1) circle (2pt);

  \draw[densityred, very thick] (6.8, 1.0) -- (9.7, 1.0) node[midway, above, font=\scriptsize, densityred] {Customer 5 ($W_5$)};
  \fill[densityred] (6.8,1.0) circle (2pt);
  \fill[densityred] (9.7,1.0) circle (2pt);

  \draw[dashed, trapred, line width=1.1pt] (4.7, 0) -- (4.7, 3.8);
  \node[above, trapred, font=\footnotesize, align=center] at (4.7, 3.8) {At time $t$: $X_t = 3$ present};
  \fill[trapred] (4.7, 1.5) circle (2.5pt);
  \fill[trapred] (4.7, 2.3) circle (2.5pt);
  \fill[trapred] (4.7, 3.1) circle (2.5pt);

  \node[anchor=north, font=\footnotesize, draw=gray!40, fill=codebg, rounded corners, inner sep=5pt] at (5.2, -0.55) {
    $\text{Total Area} = L \cdot T = (\lambda T) \cdot W \implies L = \lambda W$
  };
\end{tikzpicture}
\caption{The area geometry behind Little's law $L = \lambda W$. Each customer who visits the system occupies a horizontal interval of length $W_i$ (their sojourn time). Slicing the area accumulated under the occupancy curve vertically at each instant $t$ tallies the queue length $X_t$, integrating over $[0,T]$ to $L \cdot T$. Slicing the identical area horizontally sums the stay durations $W_i$ across all $A(T) \approx \lambda T$ arrivals, totaling $\lambda T \cdot W$. Equating the two views of the same area gives $L \cdot T = \lambda T \cdot W$, which cancels $T$ to yield $L = \lambda W$.}
\label{ch27:fig-little-area}
\end{figure}

\begin{example}[Little's law as a sanity check on M/M/1]
\label{ch27:ex-little-check}
We claimed $L = \rho/(1-\rho)$ and $W = 1/(\mu-\lambda)$. Do they obey $L = \lambda W$? Compute
\[
  \lambda W = \frac{\lambda}{\mu - \lambda}
            = \frac{\lambda/\mu}{1 - \lambda/\mu}
            = \frac{\rho}{1-\rho} = L. \checkmark
\]
Indeed Little's law is \emph{how} one usually derives $W$: find $L$ from the steady-state
distribution (easy, it is a geometric mean), then divide by $\lambda$. The same trick applied to the
queue-only quantities gives $L_q = \lambda W_q$.
\end{example}

\begin{applicationbox}[Little's law in the wild]
Because it needs so few assumptions, Little's law is a Swiss army knife far beyond telephones.
\begin{itemize}
  \item \textbf{Software / Kanban.} If a team finishes $\lambda = 5$ features per week and on average
  $L = 20$ features are ``in progress,'' then each feature takes $W = L/\lambda = 4$ weeks to ship.
  To ship faster, either finish faster or \emph{limit work in progress} --- capping $L$ directly caps
  $W$. This is the quantitative core of lean and agile flow management.
  \item \textbf{Hospitals.} With an average of $L = 60$ patients in a ward and admissions at
  $\lambda = 12$ per day, the average length of stay is $W = 5$ days. Planners invert it: target a
  stay $W$, multiply by admission rate, get the beds needed.
  \item \textbf{Inventory and supply chains.} ``Average inventory equals throughput times flow
  time.'' A warehouse holding $L$ units that move at $\lambda$ units/day implies each unit sits
  $W = L/\lambda$ days --- the basis of just-in-time analysis.
  \item \textbf{Web services.} If a service handles $\lambda = 2000$ requests/second and each spends
  $W = 50$ ms $= 0.05$ s in flight, then on average $L = 100$ requests are being processed at once
  --- exactly the concurrency a thread pool or connection limit must support.
\end{itemize}
\end{applicationbox}

\section{Beyond One Server: Networks and Variations}
\label{ch27:networks}

The single server is only a starting point. The same framework scales up, and its generalizations
form the core of applied queueing analysis.

\subsection{The M/M/$c$ queue: many servers}

With $c$ identical servers (think $c$ checkout lanes or $c$ CPU cores), the death rate depends on how
many are busy: with $n \le c$ customers, all $n$ are in service and departures fire at rate $n\mu$;
with $n > c$, all servers are busy and the rate saturates at $c\mu$. This is still a birth-death
process,
\[
  \lambda_n = \lambda, \qquad
  \mu_n = \begin{cases} n\mu, & n \le c, \\ c\mu, & n > c, \end{cases}
\]
so the product formula applies verbatim --- you just plug in the state-dependent death rates. The
result is the celebrated \emph{Erlang C} formula for the probability an arriving customer has to
wait. It is what call centers use to decide how many agents to staff, and it reveals a comforting
\emph{economy of scale}.

\begin{aestheticbox}
There is a surprising fact here, sometimes called the resource-pooling or economy-of-scale
principle. One fast server of rate $c\mu$ and $c$ slow servers each of rate $\mu$ have the same total
capacity. But \emph{a single shared line feeding many servers beats many separate lines.} Picture a
bank with one snaking queue feeding five tellers versus five separate queues. The single queue is
better, because no teller ever sits idle while someone waits in another line --- the dreaded
``I always pick the slow lane'' problem vanishes. Quantitatively, pooling $c$ servers behind one
queue lets you run each at higher utilization for the same delay, so larger systems are
\emph{relatively} more efficient. This is why cloud providers, with millions of pooled cores, can
offer low latency at high utilization that no small data center can match. The mathematics of $Q$ and
the product formula is what turns this intuition into a staffing number.
\end{aestheticbox}

\subsection{Jackson networks}

Real systems are not one queue but a web of them: a packet visits a router, then a link, then another
router; a manufactured part visits station after station. A \emph{Jackson network} is a collection of
M/M/1-style queues where, on finishing service at one node, a customer is routed to another node
(or leaves) by fixed probabilities. The astonishing \emph{product-form} theorem of Jackson (1957)
says that in steady state the nodes behave \emph{as if independent}: the joint distribution of all
the queue lengths factors,
\[
  \P(N_1 = n_1, \dots, N_k = n_k) = \prod_{i=1}^{k}\,(1-\rho_i)\rho_i^{\,n_i},
\]
each node a geometric in its own effective utilization $\rho_i$, even though the nodes clearly
interact through the routing. That factorization is what makes large networks analytically tractable
at all, and it underlies performance models of the Internet, of computer architectures, and of
factory floors.

\begin{applicationbox}[Where these chains run today]
\begin{itemize}
  \item \textbf{Computer networks.} Packet buffers in routers are queues; the M/M/1 and its kin
  predict latency and loss, and inform buffer sizing and congestion control.
  \item \textbf{Cloud and data centers.} Autoscalers watch utilization $\rho$ and the queue-length
  tail to add servers \emph{before} the $\rho/(1-\rho)$ cliff is reached. Load balancers are, quite
  literally, trying to implement single-queue pooling across many machines.
  \item \textbf{Reliability engineering.} The two-state up/down machine of
  Example~\ref{ch27:ex-twostate} generalizes to CTMC models of complex systems, computing
  availability and mean-time-to-failure for everything from aircraft to power grids.
  \item \textbf{Biology and chemistry.} Population sizes, gene expression bursts, and chemical
  reaction networks are CTMCs; the generator $Q$ is the \emph{chemical master equation} of
  stochastic kinetics.
  \item \textbf{Epidemiology.} The number of infected individuals in an outbreak is a birth-death
  process (infections are births, recoveries are deaths); the condition $\rho < 1$ becomes the
  epidemic threshold $R_0 < 1$ that separates fizzle from outbreak.
\end{itemize}
\end{applicationbox}

\section{Steady State, Stationarity, and the Balance Equations}
\label{ch27:steady}

We used the stationary distribution freely above; let us tie it down. In a discrete chain the
stationary distribution solves $\pi = \pi P$. The continuous-time analogue replaces the
one-step matrix $P$ with the generator $Q$, and the fixed-point equation becomes a flow-balance
equation.

\begin{definition}[Stationary distribution of a CTMC]
\label{ch27:def-stationary}
A probability vector $\pi$ (with $\pi_i \ge 0$, $\sum_i \pi_i = 1$) is \emph{stationary} for a CTMC
with generator $Q$ if
\[
  \boxed{\,\pi Q = 0.\,}
\]
Component-wise, for each state $j$, $\sum_{i} \pi_i q_{ij} = 0$, which rearranges to the
\emph{global balance} equation
\[
  \underbrace{\pi_j \sum_{i \ne j} q_{ji}}_{\text{rate of probability leaving } j}
  \;=\;
  \underbrace{\sum_{i \ne j} \pi_i\, q_{ij}}_{\text{rate of probability entering } j}.
\]
\end{definition}

\begin{whybox}
Why is $\pi Q = 0$ the right equation? Recall $P(t) = e^{Qt}$ and that the distribution at time $t$
starting from $\pi$ is $\pi P(t)$. ``Stationary'' means this does not change with $t$: $\pi P(t) =
\pi$ for all $t$. Differentiate at $t = 0$: $\frac{\dd}{\dd t}\big|_{0}\,\pi P(t) = \pi P'(0) = \pi
Q$, and since the left side is the derivative of a constant, it is zero. Hence $\pi Q = 0$. The
balance reading is even more intuitive: in equilibrium, the probability fluid filling each state
must drain out exactly as fast as it pours in, or the state's probability would drift. $\pi Q = 0$ is
the conservation law that says ``nothing is accumulating anywhere.''
\end{whybox}

\begin{focusbox}
Two equation levels, two strengths.
\begin{itemize}
  \item \textbf{Global balance} ($\pi Q = 0$): \emph{always} characterizes the stationary
  distribution. Total flow into each state equals total flow out. Solve this system and you have
  $\pi$.
  \item \textbf{Detailed balance} ($\pi_i q_{ij} = \pi_j q_{ji}$ for every pair): a \emph{stronger}
  condition --- flow balances across each individual edge, not just each state. When it holds the
  chain is called \emph{reversible}, and $\pi$ is far easier to find (telescoping products, as in the
  birth-death case). Detailed balance implies global balance, but not conversely.
\end{itemize}
All birth-death processes are reversible, which is exactly why the M/M/1 steady state fell out so
cleanly. When you can use detailed balance, do --- it turns a linear system into a simple recursion.
\end{focusbox}

\begin{example}[Solving a three-state CTMC by global balance]
\label{ch27:ex-threestate}
Consider states $\{A, B, C\}$ with generator
\[
  Q = \begin{pmatrix} -3 & 2 & 1 \\ 1 & -2 & 1 \\ 1 & 1 & -2 \end{pmatrix}
  \qquad(\text{rows sum to } 0).
\]
The balance equations $\pi Q = 0$ read, column by column:
\[
  -3\pi_A + \pi_B + \pi_C = 0, \qquad
  2\pi_A - 2\pi_B + \pi_C = 0, \qquad
  \pi_A + \pi_B - 2\pi_C = 0.
\]
From the first, $\pi_B + \pi_C = 3\pi_A$. Combined with $\pi_A + \pi_B + \pi_C = 1$ this gives
$4\pi_A = 1$, so $\pi_A = 1/4$. Now use $\pi_B + \pi_C = 3/4$ in the second equation,
$2\pi_A - 2\pi_B + \pi_C = 0$, i.e.\ $\tfrac12 - 2\pi_B + \pi_C = 0$. Substituting
$\pi_C = \tfrac34 - \pi_B$ gives $\tfrac12 - 2\pi_B + \tfrac34 - \pi_B = 0$, so $3\pi_B = \tfrac54$ and
$\pi_B = 5/12$, leaving $\pi_C = 3/4 - 5/12 = 1/3$. (Check: $\pi_A + \pi_B + \pi_C = \tfrac14 +
\tfrac{5}{12} + \tfrac13 = \tfrac{3}{12}+\tfrac{5}{12}+\tfrac{4}{12} = 1$, and the third equation
$\pi_A + \pi_B - 2\pi_C = \tfrac14 + \tfrac{5}{12} - \tfrac23 = 0$ holds.) The chain spends only a
quarter of its time in $A$ --- the state with the fastest exit rate ($q_A = 3$, versus $q_B = q_C =
2$) --- and more of its time in the slower-to-leave states $B$ and $C$. Fast-to-leave states shed
probability; slow-to-leave states accumulate it. (Among $B$ and $C$, which exit at the same rate, the
tie is broken by where the inflow comes from --- here $B$ edges out $C$.)
\end{example}

\begin{recoverybox}
Let us collect the through-line before the exercises. (1) A CTMC is an ordinary jump chain glued to
exponential holding times --- ``wait $\Expo(q_i)$, then jump by $P_{ij} = q_{ij}/q_i$.'' (2) The
generator $Q$ packs every rate into one matrix with zero row sums; transition probabilities are
$P(t) = e^{Qt}$, and the stationary law solves $\pi Q = 0$. (3) Birth-death processes move only to
neighbors, so detailed balance gives the product formula for $\pi$. (4) The M/M/1 queue is the
constant-rate birth-death process; its steady state is geometric, $\pi_n = (1-\rho)\rho^n$, and its
mean occupancy $L = \rho/(1-\rho)$ explodes as $\rho \to 1$. (5) Little's law $L = \lambda W$ ties
occupancy to delay by pure accounting, needing essentially no assumptions. Everything else in
queueing --- many servers, networks, priorities --- is built on these five beams.
\end{recoverybox}

\subsection{Looking ahead: zero drift and martingales}
\label{ch27:looking-ahead}

Throughout this chapter, our systems settled into equilibrium because departures outpaced arrivals:
in the M/M/1 queue, having $\rho = \lambda/\mu < 1$ meant the downward pull of service exceeded the
upward push of arrivals, gently tugging the queue backlog back toward zero whenever a burst formed.
In the language of random walks from the previous chapter, the queue length was a biased random walk
with negative drift.

What happens at the critical knife-edge boundary where arrivals and departures match exactly,
$\lambda = \mu$? At that point, the system has \textbf{zero drift}. Given the current number of
customers in the queue, its expected position a moment later is neither higher nor lower --- it is
perfectly flat. In Chapter~\ref{ch26:srw}, we met this same balance in the symmetric simple random
walk: a gambler playing a completely fair coin-toss game neither wins nor loses on average from one
round to the next.

Processes with this fair, zero-drift property are called \textbf{martingales} (named after an old
French doubling-bet system). In the next chapter, we will leave waiting lines behind and study the
pure mathematics of fair games. Martingales will give us a strikingly powerful engine: the optional
stopping theorem, which solves gambler's ruin, hitting times, and exit probabilities across complex
random processes in a single line of algebra.

\section{Exercises}
\label{ch27:exercises}

\begin{exercise}[Mental check: scaling load and service]
\label{ch27:exr-mental-scale}
An M/M/1 queue has arrival rate $\lambda$ and service rate $\mu$ with $\rho = \lambda/\mu < 1$.
Suppose customer arrivals double to $2\lambda$, and to keep up, the server's processing speed is
doubled to $2\mu$.
\begin{enumerate}[label=(\alph*)]
  \item What happens to the utilization $\rho$?
  \item What happens to the average number of customers in the system $L$?
  \item What happens to the average time a customer spends in the system $W$? Explain in one sentence
  why the customer's wait changes even though the line length distribution is identical.
\end{enumerate}
\end{exercise}

\begin{exercise}[Reading a generator]
\label{ch27:exr-readgen}
A CTMC on states $\{1,2,3\}$ has $q_{12}=2,\ q_{13}=1,\ q_{21}=4,\ q_{23}=0,\ q_{31}=1,\ q_{32}=3$
(rates per hour). (a) Write the generator matrix $Q$, filling in the diagonal. (b) What is the mean
holding time in each state? (c) Starting in state $1$, when it leaves, what is the probability it
goes to state $3$?
\end{exercise}

\begin{exercise}[Two-state availability]
\label{ch27:exr-availability}
A server fails at rate $\lambda = 1$ per week and is repaired at rate $\mu = 6$ per week. (a) Using
Example~\ref{ch27:ex-twostate}, find the long-run fraction of time the server is up. (b) What is the
mean uptime between failures, and the mean repair time? (c) If a second, independent identical server
is added and the service is ``up'' whenever \emph{at least one} server works, set up (do not fully
solve) the CTMC for the number of working servers and identify it as a birth-death process.
\end{exercise}

\begin{exercise}[M/M/1 numbers]
\label{ch27:exr-mm1numbers}
Customers arrive at a single-teller bank as a Poisson process at $\lambda = 30$ per hour; service is
exponential at $\mu = 40$ per hour. Compute (a) the utilization $\rho$; (b) the probability the
teller is idle; (c) the mean number in the system $L$; (d) the mean time in the system $W$ in
minutes; (e) the probability that there are at least $3$ customers in the system.
\end{exercise}

\begin{exercise}[The cliff]
\label{ch27:exr-cliff}
For an M/M/1 queue, the mean number in system is $L(\rho) = \rho/(1-\rho)$. (a) Show that $L$ doubles
when you go from $\rho = 2/3$ to some $\rho^*$; find $\rho^*$. (b) Compute the derivative
$\dd L/\dd\rho$ and explain in one sentence why it is so large near $\rho = 1$. (c) An operator runs
at $\rho = 0.95$ and wants to cut $L$ in half. By what factor must service rate $\mu$ increase if
$\lambda$ is fixed?
\end{exercise}

\begin{exercise}[Little's law everywhere]
\label{ch27:exr-little}
(a) A coffee shop serves $\lambda = 120$ customers per hour, and a customer spends on average
$W = 6$ minutes inside. How many customers, on average, are inside at once? (b) A development team
keeps $L = 15$ tickets in progress and closes $\lambda = 3$ per day. What is the average cycle time
$W$? (c) The team wants to cut cycle time to $3$ days without working faster (same $\lambda$). What
must $L$ become, and what policy achieves it?
\end{exercise}

\begin{exercise}[Derive the geometric steady state]
\label{ch27:exr-derive-geom}
Starting from the detailed-balance equations $\pi_n \lambda = \pi_{n+1}\mu$ for the M/M/1 queue,
(a) show by induction that $\pi_n = \rho^n \pi_0$ with $\rho = \lambda/\mu$; (b) use
$\sum_{n\ge0}\pi_n = 1$ to show $\pi_0 = 1 - \rho$, stating clearly where you need $\rho < 1$;
(c) explain in words what goes wrong physically when $\rho \ge 1$.
\end{exercise}

\begin{exercise}[Global balance by hand]
\label{ch27:exr-global}
For the three-state generator in Example~\ref{ch27:ex-threestate}, verify the claimed stationary
distribution $\pi = (1/4,\,5/12,\,1/3)$ by checking that $\pi Q = 0$ directly (compute each of the
three components and confirm they are zero).
\end{exercise}

\begin{exercise}[Mean hitting time, finite chain]
\label{ch27:exr-hitting}
A particle moves on states $\{0,1,2\}$ as a birth-death chain with birth rate $\lambda = 1$ and death
rate $\mu = 2$ (so $q_{01}=1$, $q_{10}=2$, $q_{12}=1$, $q_{21}=2$). Let $h_i$ be the expected time to
first reach state $0$ starting from state $i$. Using the first-step idea ``mean holding time in $i$
plus the expected remaining time from where you jump,'' set up equations for $h_1$ and $h_2$ and
solve them. (Hint: $h_0 = 0$; from state $i$ you wait $1/q_i$ then move.)
\end{exercise}

\begin{exercise}[Proof: detailed balance implies global balance]
\label{ch27:exr-proof-balance}
Suppose $\pi$ satisfies detailed balance for a CTMC: $\pi_i q_{ij} = \pi_j q_{ji}$ for all
$i \ne j$. Prove that $\pi$ then satisfies global balance, i.e.\ $\pi Q = 0$. (Hint: fix a state $j$,
sum the detailed-balance equations over all $i \ne j$, and recall that $q_{jj} = -\sum_{i\ne j}
q_{ji}$.) Then give a one-line argument for why the converse can fail, citing a chain that is not
reversible.
\end{exercise}

\begin{exercise}[Conceptual: why exponential?]
\label{ch27:exr-why-exp}
Explain, in a short paragraph, why a continuous-time \emph{Markov} chain is forced to have
exponentially distributed holding times. Your answer should connect the memorylessness of the
exponential distribution to the defining property that the future depends only on the present state,
not on how long the chain has already occupied it.
\end{exercise}

\begin{exercise}[Economy of scale, numerically]
\label{ch27:exr-pooling}
A service must handle $\lambda = 8$ jobs per second. Option A: one fast server at $\mu = 10$/s.
Option B: split the arrivals randomly into two streams of $4$/s each, served by two \emph{separate}
M/M/1 servers at $\mu = 5$/s each. (a) Compute $\rho$, $L$, and $W$ for Option A. (b) Compute $L$ and
$W$ \emph{per stream} for Option B and the total expected number in the whole system. (c) Which
option gives shorter average wait, and which idea from Section~\ref{ch27:networks} does this
illustrate?
\end{exercise}
